\chapter{Installation and requirements}
\label{ch:install}

\section{MCAMS inside ToScaNA}

Since November 2026 MCAMS is part of ToScaNA, as the subpackage
\code{toscana.corrections.mcams}. The code is unchanged except that the modules
import each other as a package (\code{from . import beam} instead of
\code{import beam}), and that the default source of tabulated masses and cross
sections is ToScaNA's isotope table (\code{nuclear_data.ToscanaData}), so that
\code{periodictable} is no longer needed. Make ToScaNA importable (ToScaNA
manual, chapter \emph{Installation}) and import the modules as
\begin{lstlisting}[style=shell]
from toscana.corrections.mcams.cylinder import create_cylindrical_sample
from toscana.corrections.mcams import beam, slab
\end{lstlisting}
All the examples of this manual are written this way.

\section{Requirements}

\begin{itemize}
  \item \textbf{Python} 3.10 or newer (as ToScaNA).
  \item \textbf{NumPy} (tested with 2.4) for all the calculations.
  \item \textbf{Matplotlib} (tested with 3.10) for the plots. It is only needed by the \code{plot_...}
        functions and by the demonstrations at the end of the modules.
  \item \textbf{periodictable} (optional): an alternative source of tabulated cross sections and masses,
        used only with \code{data=AtomicData()} (Chapter~\ref{ch:builder}).
\end{itemize}

\section{Files and folders}

\begin{lstlisting}[style=shell]
toscana/corrections/mcams/
    beam.py  sample_base.py  slab.py  disc.py  cylinder.py  sphere.py  sample_factory.py
    transmission_vs_angle.py  multiple_scattering.py  transmission_by_layer.py
    direct_attenuation.py  formula.py  nuclear_data.py  layer_builder.py  run_simulation.py
toscana/tests/mcams_checks.py   (self-test, run by tests/test_mcams.py)
toscana/docs/mcams/manual/      (this manual)
\end{lstlisting}

\section{Checking the installation}

Every module has a short demonstration at its end:
\begin{lstlisting}[style=shell]
python -m toscana.corrections.mcams.slab          # flat cell, partially filled
python -m toscana.corrections.mcams.cylinder      # can, partially filled
python -m toscana.corrections.mcams.disc          # round cell; beam larger than the disc
python -m toscana.corrections.mcams.sphere        # spherical container half filled
python -m toscana.corrections.mcams.formula       # chemical formulas
python -m toscana.corrections.mcams.layer_builder # layers from formulas
\end{lstlisting}
The self-test of the geometry compares the paths with a brute-force calculation and the interaction rates
with exact results; from the ToScaNA folder:
\begin{lstlisting}[style=shell]
python tests/mcams_checks.py
...
ALL PASSED
\end{lstlisting}

\section{Use in a Jupyter notebook}

Two points deserve attention in notebooks:
\begin{enumerate}
  \item Python imports a module only once per session. If you replace a file by a newer version, the notebook
        keeps using the old one until the \textbf{kernel is restarted} (or the module is reloaded with
        \code{importlib.reload}). A message such as \code{TypeError: create_slab_sample() got an unexpected keyword argument ...}
        usually means that an old version is still loaded (see Appendix~\ref{app:trouble}).
  \item Cells of a notebook can be run in any order. After changing the sample or the beam, run again the cells that depend on them.
\end{enumerate}

\section{Reproducibility}

All the functions that use random numbers accept the argument \code{rng}, which can be an integer
(a seed) or a NumPy generator created with \code{np.random.default_rng(seed)}. With the same seed the
results are identical from one run to the next. Without it, every run gives slightly different
results, as expected from a Monte Carlo method.
